\documentclass[10pt]{article}
\usepackage{tmlr}
% \usepackage[preprint]{tmlr}

\usepackage{hyperref}
\usepackage{url}
\usepackage{amsmath,amssymb,mathtools}
\usepackage{booktabs}
\usepackage{array}
\usepackage{graphicx}
\usepackage{amsthm}
\newtheorem{proposition}{Proposition}

\newcommand{\R}{\mathbb{R}}
\newcommand{\jvp}{\textsc{jvp}}
\newcommand{\fd}{\textsc{fd}}
\newcommand{\eps}{\varepsilon}

\title{An Exact, Step-Free Instrument for Second-Order Generator Geometry}

\author{\name Anonymous}

\def\month{06}
\def\year{2026}
\def\openreview{\url{https://openreview.net/forum?id=XXXXXXXXXX}}

\begin{document}
\maketitle

\begin{abstract}
Second-order generator geometry is in live use: geodesic interpolation, curvilinear
edits, and Hessian disentanglement priors all reach for the curvature of the generated
manifold. Exact higher-order automatic differentiation is not new, but we use it as a
step-free \emph{measurement referee} and find that the finite-difference estimator a
practitioner reaches for -- a central second difference of the generator -- is step-fragile.
On our StyleGAN sweep it exhibits an empirical step-selection trilemma: the steps that best
preserve magnitude accuracy, signed bias, and rank fidelity are different, with a $\geq45\%$
magnitude-error floor of order $\sqrt{\eps_{\mathrm{mach}}}$ at the best shared step
(Proposition~\ref{prop:floor}) that fp64 does not remove, while the analogous first-order
quantity has a common stable step region. We give an exact instrument that returns the
second-order pushforward to machine precision: composed forward-mode automatic
differentiation, with no step, to any order, in memory a small constant multiple of one
forward pass. We validate it two ways -- a first-order finite-difference control agrees
with it to $2.4\%$, and on an analytic referee it reaches $2.5\times10^{-17}$ against the
best finite difference's $8.6\times10^{-9}$. The decision this forces is deterministic:
to probe second-order generator geometry, use the exact instrument, not finite
differences; any method whose per-sample estimate is a finite-difference generator Hessian
-- the Hessian Penalty among them -- is exposed to this step-selection ambiguity (averaging
or stochastic objectives may temper it; \S\ref{sec:scope}). As a case study, the exact
instrument audits a subtle, anchor-dependent second-order signal: the curvature's partial
correlation with edit nonlinearity is significant in aggregate (pooled permutation
$p\approx0.002$) but does not transfer across latent anchors (cross-anchor $95\%$ CI
$[-0.09,+0.43]$ spans zero), so first-order displacement remains the safe default. Because a
fixed-step finite difference cannot be trusted for scale, sign, and ranking at once, only the
exact reference can establish which apparent findings are anchor-stable.
\end{abstract}

\section{Introduction}
\label{sec:intro}
Editing methods for pretrained GANs increasingly treat the latent space as a Riemannian
manifold and reach for its second-order structure. Geodesic interpolation replaces linear
interpolation using a pullback metric \citep{arvanitidis2018,shao2018,chen2018}; curvilinear
edit paths are constructed to be nonlinear by design \citep{aoshima2023}; and second-order
Hessian information is used as a disentanglement prior \citep{peebles2020}. Each of these
needs the curvature of the generator -- the change of the Jacobian along an edit -- and the
estimator a practitioner reaches for is a finite-difference second difference of the
generator.

That estimator has no defensible setting. A curvature estimate must get three things right -- its
magnitude, its sign, and the ranking of directions it induces -- and for the second-order difference on
StyleGAN these three are optimized at \emph{different} steps: no single step preserves all of them, and
the signed error even changes sign across the usable range (Section~\ref{sec:real},
Table~\ref{tab:trilemma}). The first-order difference of the same code has no such conflict -- one step
region is at once accurate, unbiased, and rank-faithful -- which is why first-order Jacobians and pullback
metrics have been safe to finite-difference while second-order curvature is not. The magnitude axis of
this conflict is an irreducible relative-error floor of order $\sqrt{\eps_{\mathrm{mach}}}$
(Proposition~\ref{prop:floor}): on StyleGAN no shared step beats a $45\%$ mean error ($38\%$ median),
identical in fp32 and fp64, so it is the generator's nonlinearity, not roundoff. With no step a
practitioner can defensibly choose, finite differences do not yield a curvature one can trust for
magnitude, sign, and ranking at once.

The fix is exact rather than better-tuned. The second-order term is the directional
derivative of the first, so composing forward-mode automatic differentiation with itself
returns the second-order pushforward to machine precision -- no step, any order, in memory
a small constant multiple of one forward pass (Section~\ref{sec:inst}). We validate it
against a first-order finite-difference control that agrees to $2.4\%$ and against an
analytic referee where it reaches $2.5\times10^{-17}$ while the best finite difference
manages only $8.6\times10^{-9}$.

This settles a decision that the literature has left implicit: to probe second-order
generator geometry, use the exact instrument, not finite differences. With the exact
quantity in hand we then revisit a concrete question a practitioner faces -- does the
curvature of an edit direction predict how \emph{nonlinear} the edit is, how far the
rendered path bows away from the straight chord? -- as a case study that shows what the
instrument buys.

Our contribution is not a new higher-order AD algorithm; it is a measurement study for
second-order generator geometry. We instantiate exact composed-\jvp{}s as a deterministic
per-direction referee, show that finite differences face an empirical step-selection trilemma
on StyleGAN, and use the exact referee to audit a curvature--nonlinearity finding that appears
significant at one anchor but is anchor-fragile under replication. Concretely:

\begin{itemize}
\item A non-recoverability result (Proposition~\ref{prop:floor}) for the magnitude of the
central second difference -- a relative-error floor of order $\sqrt{\eps_{\mathrm{mach}}}$,
$\geq45\%$ on StyleGAN and identical in fp32/fp64 -- and its empirical extension to a
step-selection \emph{trilemma}: magnitude, signed bias, and rank fidelity are optimized at
different steps (Table~\ref{tab:trilemma}), whereas the first difference of the same code is
accurate, unbiased, and rank-faithful together near its optimal step ($2.4\%$)
(Sections~\ref{sec:fd},~\ref{sec:real}).
\item An exact, step-free instrument -- composed forward-mode AD -- that returns the
second-order pushforward to machine precision in memory a small constant multiple of one
forward pass, giving the $k$-th-order pushforward from $k$ nested sweeps
(Section~\ref{sec:inst}).
\item A deterministic prescription: to probe second-order generator geometry, use the exact
instrument, not finite differences; any method built on a finite-difference generator
Hessian -- the Hessian Penalty \citep{peebles2020} -- inherits the $\geq45\%$ floor
(Section~\ref{sec:scope}).
\item A case study the instrument makes answerable (Section~\ref{sec:negative}): with the
estimator confound removed, the exact second-order signal for edit nonlinearity is real but
anchor-dependent -- significant on two of five anchors and in aggregate (pooled permutation
$p\approx0.002$), yet with a cross-anchor estimate of $+0.17$ whose $95\%$ CI spans zero, so
it does not transfer and first-order displacement is the safe default. At partial
correlations of $0.1$--$0.37$ this signal sits below the finite-difference floor and is
measurable only with the exact instrument.
\end{itemize}

The instrument is the airtight contribution; Proposition~\ref{prop:floor} is proved and the
exact second-order pushforward is validated to machine precision. The case study is reported
at the strength its evidence supports: a genuine but anchor-dependent second-order signal
that only the exact instrument can resolve.

\section{Setup}
\label{sec:setup}
Let $G:\R^{d}\to\R^{3HW}$ be a pretrained generator and $z$ a latent code. For a unit
direction $v$ the edit path is $\alpha\mapsto G(z+\alpha v)$. We compare three quantities along
this path at $\alpha=0$.

\paragraph{Prediction target (estimator-free).} The \emph{path nonlinearity} $\mathrm{NL}(z,v)$
is how far the rendered path departs from the straight chord between its endpoints, measured by
rendering only: sample $K$ points $\alpha_k\in[-T,T]$, form the chord between $G(z-Tv)$ and
$G(z+Tv)$, and average the normalized image-space deviation of the interior points from the
chord. This is exactly ``how wrong is linear interpolation'' over the edit range, and it uses
no derivative estimator, so it cannot favor either proxy.

\paragraph{First-order proxy (baseline).} The displacement $\lVert J_z v\rVert =
\lVert\partial_\alpha G\rVert$, the exact first-order \jvp{}. This is the cheap quantity a
practitioner already has.

\paragraph{Second-order proxy.} The curvature magnitude $\lVert\partial_\alpha^2 G\rVert$.
Whether this predicts $\mathrm{NL}$ beyond the first-order proxy is the case study we return
to in Section~\ref{sec:negative}; obtaining it reliably is the problem the instrument solves.

\section{Second differences cannot recover generator curvature}
\label{sec:fd}
The obvious way to obtain $\partial_\alpha^2 G$ without automatic differentiation is a central
second difference. It does not work in the precision generators run in.

\begin{proposition}[Second-order non-recoverability]
\label{prop:floor}
Let $f\in C^4$ near $\alpha$ with $f^{(4)}$ bounded and $f''(\alpha)\neq0$, evaluated in
floating point with unit roundoff $\eps_{\mathrm{mach}}$ under a standard rounding model
\citep{higham2002}. The central second difference $D^2_\eps f=\big(f(\alpha+\eps)-2f(\alpha)+
f(\alpha-\eps)\big)/\eps^2$ has relative error
\[
\frac{|D^2_\eps f-f''(\alpha)|}{|f''(\alpha)|}
\;\lesssim\; C_T\,\eps^2 + C_R\,\eps_{\mathrm{mach}}/\eps^2,
\qquad C_T=\frac{|f^{(4)}(\alpha)|}{12\,|f''(\alpha)|},\;\;
C_R=O\!\Big(\frac{|f(\alpha)|}{|f''(\alpha)|}\Big),
\]
whose minimum over the step is of order $\sqrt{C_T C_R\,\eps_{\mathrm{mach}}}$, attained at
$\eps^\star\sim\eps_{\mathrm{mach}}^{1/4}$; no step does better. The central first difference, by
contrast, attains a floor of order $\eps_{\mathrm{mach}}^{2/3}$. For a vector map
$G:\R^{d}\to\R^{3HW}$ the bound holds componentwise, so a single global step is governed by the
worst component.
\end{proposition}

The truncation term $\tfrac{\eps^2}{12}f^{(4)}$ and the cancellation term
$O(\eps_{\mathrm{mach}}|f|/\eps^2)$ cannot be small together; the exponent $\tfrac12$ is intrinsic
to the operation and the constant carries the generator's nonlinearity $|f^{(4)}f|/f''^2$, which is
large for a deep network (derivation in Appendix~\ref{app:deriv}). In double precision the floor is
about $10^{-8}$ in the benign case; in the single precision generators use
($\eps_{\mathrm{mach}}\approx1.2\times10^{-7}$) it is about $3\times10^{-4}$ benign and far worse for
a strongly nonlinear $G$, where Section~\ref{sec:real} measures $\geq45\%$. The first difference is
different: a milder floor $\eps_{\mathrm{mach}}^{2/3}$ with a wide usable window, so first-order
Jacobians are safe to finite-difference while second-order curvature is not. This asymmetry is also
our control: wherever the first-order finite difference agrees with the exact \jvp{}, the pipeline is
correct and any second-order disagreement is real.

To check the constants against ground truth, let $f(\alpha)=[\,e^{0.7\alpha},\sin 3\alpha,
(\alpha+0.4)^3,\tanh 2\alpha,\alpha^2\cos\alpha\,]$, with derivatives known in closed form. At
$\alpha=0.3$ in double precision (Table~\ref{tab:toy}) the central second-order relative error is
U-shaped, bottoming at $8.6\times10^{-9}$ near $\eps=10^{-4}$ and diverging to $0.29$ by
$\eps=10^{-8}$; the composed \jvp{} reaches $2.5\times10^{-17}$, so the best finite difference is
$3.5\times10^{8}$ times worse, matching the Proposition~\ref{prop:floor} bound with the $1/12$
constant.

\begin{table}[h]\centering
\begin{tabular}{rcl}
\toprule
$\eps$ & central \fd{} 2nd-order rel.\ err. & \\
\midrule
$10^{-1}$ & $7.9\times10^{-3}$ & (truncation) \\
$10^{-2}$ & $8.0\times10^{-5}$ & \\
$10^{-4}$ & $8.6\times10^{-9}$ & (best) \\
$10^{-6}$ & $3.0\times10^{-5}$ & \\
$10^{-8}$ & $2.9\times10^{-1}$ & (cancellation) \\
\midrule
composed \jvp & $2.5\times10^{-17}$ & (exact) \\
\bottomrule
\end{tabular}
\caption{Second-order accuracy on the analytic toy at $\alpha=0.3$, double precision. No
finite-difference step approaches the exact composed \jvp{}.}
\label{tab:toy}
\end{table}

\section{The non-recoverability is real on a generator}
\label{sec:real}
On StyleGAN-bedroom at $256{\times}256$ in fp32, the best single step $\delta\!\approx\!3$ shared across
edit directions leaves the central second-difference curvature with a $45\%$ mean relative-error floor
($38\%$ median across directions), rising on both sides and exploding as $\delta\to0$
(Figure~\ref{fig:floor}). On matched inputs the floor is identical in fp32 and fp64 ($47\%$ in both; the
gap from the $45\%$ best-step value reflects the input sample, not precision), confirming it is truncation
against the generator's nonlinearity rather than roundoff. This $\geq45\%$ figure is the \emph{magnitude}
axis of Proposition~\ref{prop:floor}; on a real generator the difficulty is in fact three-fold.

\paragraph{An empirical step-selection trilemma.} Magnitude accuracy is not the only thing a
second-order estimate must get right. A practitioner may need the curvature \emph{magnitude} (geodesic
energy, regularizer strength), its \emph{sign}/unbiasedness (a thresholded decision), or the
\emph{ranking} of directions by curvature (which edit is ``most curved''). Across our $96$
direction--anchor pairs these three criteria are optimized at \emph{different} steps
(Table~\ref{tab:trilemma}): magnitude error is least at $\delta=3.0$ ($45\%$), the signed bias is
smallest at $\delta=2.0$ ($+5\%$), and rank fidelity is highest at $\delta\approx0.2$ (Spearman $0.983$).
The magnitude-optimal step gives lower rank fidelity than the rank-optimal step, while the rank-optimal
step incurs a $355\%$ magnitude error; and the signed bias crosses from underestimation at large steps to
overestimation at small steps, so there is no ``safe side'' to err on. This three-way conflict is
\emph{empirical} -- Proposition~\ref{prop:floor} bounds only the magnitude axis; the bias and rank axes
are measured, not proved -- and it is specific to second order: the analogous first-order quantity is
accurate, unbiased, and rank-faithful together in a common region near $\delta=0.02$ ($2.4\%$ error,
$-0.0\%$ bias, Spearman $0.999$; right columns of Table~\ref{tab:trilemma}). The first-order central
difference of the same code thus agrees with the exact \jvp{} to $2.4\%$, so the pipeline is correct and
the second-order conflict is genuine. The exact composed-\jvp{} instrument removes the step axis
altogether and serves as the reference against which all three criteria are even defined. These specific
step values are properties of this generator's nonlinearity; what we expect to generalize is the
\emph{existence} of the conflict -- a magnitude floor coexisting with an opposite-signed cancellation
regime -- not the particular steps, which is why we state the trilemma as empirical and confined to this
StyleGAN sweep.

\begin{table}[t]\centering
\begin{tabular}{r l rrr rr}
\toprule
& & \multicolumn{3}{c}{second order (FD vs exact)} & \multicolumn{2}{c}{first order} \\
\cmidrule(lr){3-5}\cmidrule(lr){6-7}
step $\delta$ & role & mag.\ err & signed bias & Spearman & mag.\ err & Spearman \\
\midrule
$8.0$          & large step        & $74\%$           & $-74\%$            & $0.71$           & $64\%$          & $0.69$            \\
$\mathbf{3.0}$ & 2nd mag-best      & $\mathbf{45\%}$  & $-26\%$            & $0.90$           & $44\%$          & $0.87$            \\
$\mathbf{2.0}$ & 2nd bias-best     & $56\%$           & $\mathbf{+5\%}$    & $0.93$           & $37\%$          & $0.90$            \\
$\mathbf{0.2}$ & 2nd rank-best     & $355\%$          & $+350\%$           & $\mathbf{0.983}$ & $9\%$           & $0.995$           \\
$0.02$         & 1st-order stable  & $6773\%$         & $+6773\%$          & $0.72$           & $\mathbf{2.4\%}$ & $\mathbf{0.999}$ \\
\bottomrule
\end{tabular}
\caption{Step-selection trilemma for second-order generator curvature on StyleGAN-bedroom ($96$
direction--anchor pairs; central finite difference vs the exact composed-\jvp{}). The steps that best
preserve magnitude ($\delta=3.0$), signed bias ($\delta=2.0$), and rank fidelity ($\delta=0.2$) are
different, and the signed bias crosses from underestimation at large steps to overestimation at small
steps. The analogous first-order quantity (right columns) is accurate, unbiased, and rank-faithful in a
common region near $\delta=0.02$. Bold marks the per-criterion optimum. Full $16$-step sweep in
Appendix~\ref{app:trilemma}.}
\label{tab:trilemma}
\end{table}

\begin{figure}[t]\centering
\includegraphics[width=\linewidth]{fig1_floor.pdf}
\caption{Non-recoverability of Proposition~\ref{prop:floor} on StyleGAN. The first-order
finite difference (blue) has a usable minimum; the second-order difference (red, mean over directions)
never falls below $\approx45\%$ at any common step, and the floor is identical in fp32 and fp64.}
\label{fig:floor}
\end{figure}

\section{An exact, step-free instrument}
\label{sec:inst}
The second-order term is the directional derivative of the first, so composing forward-mode AD with
itself,
\[
\partial^2_\alpha G(z+\alpha v)=\partial_\alpha\big[\partial G(z+\alpha v)[v]\big][v],
\]
is the exact derivative of the implemented graph and is evaluated by nesting two \jvp{} calls (in
PyTorch, \texttt{torch.func.jvp} of a function that returns a \jvp{}) \citep{paszke2019,griewank2008}.
Exact higher-order directional derivatives by composed/forward-mode AD are long-standing
machinery -- the Hessian-vector product at gradient cost \citep{pearlmutter1994}, hyper-dual
numbers for exact second derivatives without a step \citep{fike2011}, and second-order
forward-mode AD \citep{baydin2018,cobb2024} -- so the instrument carries no algorithmic
novelty; its role is exactness at generator scale.
No Jacobian or Hessian is materialized, so memory is a small constant multiple of a single forward
evaluation -- about two for the second-order term -- rather than scaling with the output dimension, and
$k$ nested sweeps give the $k$-th-order pushforward. The instrument needs only a differentiable forward
graph and applies to any such architecture (a short per-generator adapter sufficed in our experiments). Its purpose here is to make the
prediction question of Section~\ref{sec:negative} answerable with the \emph{exact} second-order quantity,
so that whatever the answer turns out to be cannot be blamed on the estimator.

\section{A case study the instrument makes answerable}
\label{sec:negative}
We now use the exact instrument to ask a concrete question: does an edit direction's curvature
predict the estimator-free nonlinearity $\mathrm{NL}$ beyond the first-order displacement? This is
exactly the question the finite-difference estimator's $\geq45\%$ floor (Section~\ref{sec:fd})
would make uninterpretable -- any answer could be the estimator talking. The exact pushforward
removes that excuse, so the answer, whatever it is, is about the geometry. We use StyleGAN-bedroom,
$50$ edit directions (the $8$ annotated attribute boundaries plus $42$ random $w$-directions), and
average each quantity over $12$ latent anchors per direction; we repeat the entire correlation
across five independent anchor seeds and report the distribution. Rank correlations are Spearman;
partial correlations residualize ranks; confidence intervals are bootstrap over directions; the
first-order control of Section~\ref{sec:real} passes throughout ($2.4\%$).

\paragraph{A single anchor looks strong but is fragile.} On one anchor set ($N{=}24$ named
directions, a subset of the fifty), exact curvature correlates with nonlinearity at $\rho=0.47$
($p=0.012$) and, controlling for first-order displacement, the partial correlation is $+0.45$
($p=0.016$, $10^4$ permutations) -- a substantial second-order signal. It is fragile here:
displacement and curvature are nearly collinear ($\rho_{c,J}=0.93$), and the raw curvature
correlation drops to $0.31$ ($p=0.075$) once two heavy-tailed directions are removed. A single
anchor cannot decide the question.

\paragraph{Across anchors the signal is real but anchor-dependent.} On five anchors (50 directions
each, Figure~\ref{fig:neg}b) the partial correlation $\rho(\text{curv},\mathrm{NL}\mid\lVert
J_zv\rVert)$ is $+0.10,\,+0.19,\,-0.15,\,+0.37,\,+0.33$. We report both tests symmetrically. By
permutation the per-anchor $p$-values are $0.25,\,0.095,\,0.86,\,0.0038,\,0.010$, so the signal is
significant on two of five anchors, and both survivors clear a Bonferroni threshold of $.01$; a
pooled permutation test across anchors is significant (Fisher $p\approx0.002$). By bootstrap, on the
other hand, every anchor's $95\%$ CI includes zero (e.g.\ seed $34$ $[-0.076,+0.648]$, seed $35$
$[-0.085,+0.611]$). The two readings are consistent: there is a genuine aggregate effect, but it is
not pinned down at any single anchor. The cross-anchor (random-effects) estimate is $+0.17$ with
$95\%$ CI $[-0.09,+0.43]$, falling to $+0.06$ ($95\%$ CI $[-0.23,+0.34]$, median $+0.09$) once chord
length is also controlled. So the effect does not transfer reliably to a new anchor.

\begin{figure}[t]\centering
\includegraphics[width=\linewidth]{fig2_negative.pdf}
\caption{The exact second-order signal for edit nonlinearity is real but anchor-dependent.
(a) On a single anchor, first-order displacement and second-order curvature are nearly collinear
predictors of path nonlinearity. (b) Across five anchors, the partial correlation of curvature given
displacement is significant on $2/5$ by permutation (seeds $34,35$) and in aggregate, while each
anchor's bootstrap $95\%$ CI includes zero; the cross-anchor estimate $+0.17$ has a CI that spans
zero, so the effect does not transfer.}
\label{fig:neg}
\end{figure}

\paragraph{Reading the case study.} The exact second-order signal is real in aggregate -- two of
five anchors significant, pooled $p\approx0.002$ -- but anchor-dependent, and it does not transfer:
for a practitioner who cannot know in advance which anchor regime an edit sits in, first-order
displacement is the safe default. The point for this paper is methodological. The $\geq45\%$ figure is
a \emph{magnitude} error, whereas these case-study correlations are rank/residual statistics, so we do
not read the floor as directly burying a correlation of size $0.1$--$0.37$. The instrument-necessity is
subtler: a fixed-step finite difference faces the trilemma of Section~\ref{sec:real}, so its curvature
cannot be trusted for scale, sign, and ranking at once, and the practitioner cannot know in advance which
step preserves the property a given audit needs. A nuanced, anchor-dependent second-order effect is
therefore reliably auditable only with the exact reference, which removes the step choice. That is the
case for it.

\paragraph{Why the question is a priori reasonable.} The result is not forced by collinearity:
the chord length -- raw global displacement -- is uncorrelated with nonlinearity ($\rho=0.065$), so
$\mathrm{NL}$ is not merely a re-encoding of how far the edit moves. The leading non-trivial term of
the chord deviation is $\tfrac12\alpha^2\partial^2_\alpha G$, so second-order curvature is the
\emph{a priori} right-order predictor of nonlinearity; the instrument lets us test that prediction
honestly rather than through the estimator's floor, and the honest answer is the mixed,
anchor-dependent one above.

\section{Prescription, scope, and limitations}
\label{sec:scope}
\paragraph{Prescription.} To probe second-order generator geometry, use the exact instrument, not
finite differences. This is a deterministic consequence of Proposition~\ref{prop:floor}: a
finite-difference generator Hessian has a $\geq45\%$ relative-error floor in single precision that no
step size or precision upgrade removes, while the composed-\jvp{} pushforward is exact to machine
precision at the cost of a small constant multiple of one forward pass. The same floor governs any
method built on a finite-difference second derivative of a generator. The Hessian Penalty
\citep{peebles2020} estimates exactly this second difference
$(1/\eps^2)[G(z+\eps v)-2G(z)+G(z-\eps v)]$; it is a stochastic training-time regularizer, so
averaging across samples may temper the error, but each per-sample estimate it averages sits on the
floor of Proposition~\ref{prop:floor}, and the exact composed-\jvp{} Hessian-vector product is a
drop-in replacement at the same forward-pass cost.

\paragraph{A methodological takeaway.} More broadly, a second-order generator-geometry claim
should report an exact or validated estimator and replication across anchors. Both failure modes
we document can manufacture an apparent effect that does not survive a control: the
finite-difference step-selection trilemma (Section~\ref{sec:real}) lets the choice of step decide
whether a magnitude, sign, or ranking conclusion holds, and single-anchor variability
(Section~\ref{sec:negative}) lets the choice of base latent decide whether a correlation is
significant.

\paragraph{On the case study.} The second-order signal for predicting edit nonlinearity is real in
aggregate but anchor-dependent (Section~\ref{sec:negative}); it does not transfer reliably across
anchors, so for a practitioner who cannot fix the anchor regime, first-order displacement
$\lVert J_z v\rVert$ -- a single \jvp{} -- is the safe default. That the signal is auditable at all
is itself a result for the instrument: a fixed-step finite difference cannot be trusted for the scale,
sign, and ranking of curvature at once (Section~\ref{sec:real}), so a partial correlation this small
could not be reliably attributed to geometry rather than step choice without the exact reference.

\paragraph{Scope and limitations.} The case study is confined to what we measured -- one generator
family (StyleGAN-bedroom), the path-nonlinearity target, and curvature versus first-order
displacement -- and the cross-anchor effect is modest ($+0.17$, CI spanning zero), so it stands only
within this setting; we did not test other generators (StyleGAN-FFHQ, diffusion U-Nets) or other
targets (identity preservation, attribute leakage). StyleGAN-bedroom is nonetheless a fair stress
test for the central result: its strong nonlinearity is the worst case for the finite-difference
floor, so the $\geq45\%$ figure is if anything conservative for milder generators. The
non-recoverability result (Proposition~\ref{prop:floor}) is architecture-independent and the
instrument applies to any differentiable generator.

\section{Related work}
\label{sec:related}
Latent-geometry analyses of generative models use the first-order pullback metric $G=J^\top J$ to define
geodesics and argue that geodesic interpolation improves on linear interpolation
\citep{arvanitidis2018,shao2018,chen2018,yang2018}; these methods apply a curved path everywhere and rely
on the generator's higher-order structure without separating how reliably that structure can be computed.
Curvilinear editing \citep{aoshima2023} constructs nonlinear edit paths to be commutative. Second-order
information appears as a training-time disentanglement prior in the Hessian Penalty \citep{peebles2020},
which estimates the Hessian by exactly the finite-difference second difference our
Proposition~\ref{prop:floor} shows is non-recoverable. Closest to our object,
\citet{lee2023} regularize the explicit (extrinsic) curvature of a deep generative decoder
using second-order automatic differentiation, with a stochastic Hutchinson trace estimator
on low-dimensional autoencoders as a training-time prior; we instead compute the exact
per-direction second-order pushforward of pretrained production-scale generators as a
measurement instrument and quantify when the finite-difference alternative fails. Recent work computes curvature of generative
manifolds by careful grid finite differences \citep{koulakis2026}, a regime distinct from the
single-anchor second difference studied here. The finite-difference step-size tradeoff is classical
\citep{press2007,higham2002} and higher-order forward-mode AD is standard \citep{pearlmutter1994,fike2011,griewank2008,baydin2018,cobb2024}; our
contribution is to bring them together into an exact, step-free instrument for second-order generator
geometry, to establish that the finite-difference alternative is non-recoverable on real generators, and
to show on a case study that the exact reference can audit a subtle second-order signal whose scale,
sign, and ranking a fixed-step finite difference cannot be trusted to preserve at once.

\section{Conclusion}
Finite differences cannot reliably recover the curvature of a generator. In single precision the central
second difference has a $\geq45\%$ magnitude-error floor that no step or precision removes
(Proposition~\ref{prop:floor}), and on a real generator this widens into an empirical step-selection
trilemma: the steps that best preserve magnitude, signed bias, and rank fidelity are different, while the
first-order quantity stays accurate, unbiased, and rank-faithful in a common step region
(Section~\ref{sec:real}). The exact composed-\jvp{} instrument can recover
it, to machine precision and at the cost of a small constant multiple of one forward pass. The
prescription follows directly: to probe second-order generator geometry, use the exact instrument, not
finite differences -- any method built on a finite-difference generator Hessian depends on an
unvalidated step choice (averaging may temper it; \S\ref{sec:scope}). As an
illustration of the instrument in use -- not a general claim about curvature -- we audit one such
question: the curvature--nonlinearity effect is real in aggregate but anchor-dependent and does not
transfer, so first-order displacement remains the safe default for a practitioner who cannot fix the
anchor regime.
The single-anchor estimate that looked strong and did not replicate is a reminder that such a signal
needs replication across anchors before it is trusted, in the spirit of multi-seed audits of
apparent structure elsewhere \citep{locatello2019,hewitt2019}.

\section*{Reproducibility statement}
The instrument, the analytic toy, and the StyleGAN scripts (Sections~\ref{sec:fd}--\ref{sec:negative})
with fixed seeds and the \texttt{metrics.json} behind every number -- including the per-anchor prediction
tables (\texttt{prediction\_multianchor\_table.json}, \texttt{prediction\_single\_anchor\_gonogo.json})
and the step-selection trilemma of Tables~\ref{tab:trilemma} and~\ref{tab:fullsweep}, whose three axes
are regenerated from the per-record finite-difference evidence by \texttt{compute\_trilemma\_table.py}
into \texttt{trilemma\_table.json} -- are released on acceptance. The
generator is the public StyleGAN-bedroom-256 checkpoint \citep{karras2019}; all experiments run on a
single 8\,GB GPU.

\bibliography{refs}
\bibliographystyle{tmlr}

\newpage
\appendix
\section{Derivation of Proposition~\ref{prop:floor}}
\label{app:deriv}
Taylor expansion of $f(\alpha\pm\eps)$ to fourth order gives $D^2_\eps f = f''(\alpha) +
\tfrac{\eps^2}{12}f^{(4)}(\alpha)+O(\eps^4)$ for the exact arithmetic, the truncation term. In floating
point each evaluation carries error $O(\eps_{\mathrm{mach}}|f|)$; dividing the differenced numerator by
$\eps^2$ amplifies it to $O(\eps_{\mathrm{mach}}|f|/\eps^2)$, the cancellation term. Summing and
minimizing over $\eps$ gives $\eps^\star\sim(\eps_{\mathrm{mach}}C_R/C_T)^{1/4}\sim
\eps_{\mathrm{mach}}^{1/4}$ and floor $\sim\sqrt{C_TC_R\,\eps_{\mathrm{mach}}}$. The first difference
divides by $\eps$, giving cancellation $O(\eps_{\mathrm{mach}}|f|/\eps)$, optimum
$\eps_{\mathrm{mach}}^{1/3}$, floor $\eps_{\mathrm{mach}}^{2/3}$.

\section{Experimental setup}
\label{app:exp}
Path nonlinearity uses $T=3$ and $K=33$ render points. The second-order proxy is the mean absolute
$\partial^2_\alpha G$ via two nested \texttt{torch.func.jvp} calls; the first-order proxy is the \jvp{}
norm. Directions are the $8$ HiGAN attribute boundaries (placed in their default manipulation layers) and
$42$ random unit $w$-directions. The multi-anchor analysis uses anchor seeds $31$--$35$, $12$ base
latents each; permutation tests use $10^4$ resamples; bootstrap CIs use $2000$ resamples over directions.

\section{Full finite-difference step sweep}
\label{app:trilemma}
Table~\ref{tab:trilemma} in the body shows five representative steps. Table~\ref{tab:fullsweep} below
gives the complete $16$-step grid behind it, for both the second- and first-order finite differences
against the exact composed-\jvp{}, averaged over all $96$ direction--anchor pairs; the body's five rows
are a subset, not a selection. Each criterion's optimum (least magnitude error, least $|$signed bias$|$,
highest Spearman) is reached at a different second-order step, while the first-order columns are jointly
near-optimal across a broad region around $\delta=0.02$--$0.05$.

\begin{table}[h]\centering
\small
\begin{tabular}{r rrr rrr}
\toprule
& \multicolumn{3}{c}{second order (FD vs exact)} & \multicolumn{3}{c}{first order (FD vs exact)} \\
\cmidrule(lr){2-4}\cmidrule(lr){5-7}
step $\delta$ & mag.\ err & signed bias & Spearman & mag.\ err & signed bias & Spearman \\
\midrule
$8.0$   & $74.5\%$       & $-74.4\%$      & $0.712$ & $63.8\%$  & $-63.8\%$ & $0.687$ \\
$5.0$   & $57.3\%$       & $-55.6\%$      & $0.818$ & $54.0\%$  & $-54.0\%$ & $0.770$ \\
$3.0$   & $45.2\%$       & $-26.2\%$      & $0.897$ & $44.1\%$  & $-44.1\%$ & $0.867$ \\
$2.0$   & $56.1\%$       & $+4.9\%$       & $0.934$ & $37.1\%$  & $-37.1\%$ & $0.903$ \\
$1.5$   & $73.6\%$       & $+31.7\%$      & $0.950$ & $32.4\%$  & $-32.4\%$ & $0.924$ \\
$1.0$   & $106.4\%$      & $+76.1\%$      & $0.965$ & $26.2\%$  & $-26.2\%$ & $0.948$ \\
$0.7$   & $142.0\%$      & $+121.6\%$     & $0.975$ & $21.3\%$  & $-21.3\%$ & $0.967$ \\
$0.5$   & $184.2\%$      & $+170.3\%$     & $0.981$ & $17.4\%$  & $-17.4\%$ & $0.980$ \\
$0.3$   & $266.9\%$      & $+258.7\%$     & $0.983$ & $12.5\%$  & $-12.5\%$ & $0.991$ \\
$0.2$   & $355.2\%$      & $+350.4\%$     & $0.983$ & $9.4\%$   & $-9.4\%$  & $0.995$ \\
$0.1$   & $622.8\%$      & $+621.0\%$     & $0.970$ & $5.5\%$   & $-5.5\%$  & $0.999$ \\
$0.05$  & $1{,}429\%$    & $+1{,}429\%$   & $0.906$ & $3.1\%$   & $-2.9\%$  & $0.999$ \\
$0.02$  & $6{,}773\%$    & $+6{,}773\%$   & $0.721$ & $2.4\%$   & $-0.0\%$  & $0.999$ \\
$0.01$  & $25{,}023\%$   & $+25{,}023\%$  & $0.575$ & $5.7\%$   & $+4.5\%$  & $0.996$ \\
$0.005$ & $93{,}346\%$   & $+93{,}346\%$  & $0.527$ & $15.0\%$  & $+14.4\%$ & $0.986$ \\
$0.001$ & $1{,}985{,}856\%$ & $+1{,}985{,}856\%$ & $0.409$ & $139.1\%$ & $+139.0\%$ & $0.765$ \\
\bottomrule
\end{tabular}
\caption{Complete $16$-step finite-difference sweep on StyleGAN-bedroom ($96$ direction--anchor pairs),
second- and first-order, against the exact composed-\jvp{}. The second-order magnitude optimum
($\delta=3.0$), signed-bias optimum ($\delta=2.0$), and rank optimum ($\delta\approx0.2$--$0.3$) are at
distinct steps; the first-order quantity is jointly accurate, unbiased, and rank-faithful around
$\delta=0.02$--$0.05$.}
\label{tab:fullsweep}
\end{table}

\end{document}
